% Section 1 --- The problem
% the only applied change is the approved citation-key fix
% Dunne1996 -> Dunne1996Langmuir.
\section{The problem}\label{sec:problem}

A materials-discovery calculation often begins with a simple request: rank porous crystals by how strongly they adsorb a gas at low pressure. For a zeolite or a metal-organic framework, the desired quantities are usually the Henry constant, \(K_H\), and the zero-coverage isosteric heat, \(Q_{st}\). They are attractive screening targets because they are cheaper than full isotherms and because they seem to measure the first contact between a guest molecule and a host framework.

Widom insertion appears to supply exactly these quantities. In its simplest form, the method inserts a ghost molecule into a host, evaluates the host-guest energy, and averages the corresponding Boltzmann weight. In the dilute limit this is a well-defined statistical-mechanical estimator \cite{FrenkelSmit2002}. The apparent simplicity is the danger. The estimator may be exact for a specified model, but the prediction is not exact for the material unless the model is the right one.

A Widom adsorption number is therefore not a freestanding fact about a crystal. It is the result of a recipe. The recipe includes the crystallographic structure, atom typing, force field, gas model, cross-pair mixing rule, electrostatics, simulation backend, sampling rule, and experimental reference observable. Changing any one of these ingredients changes the scientific meaning of the number. A calculation can be reproducible and still be wrong if the structure is the wrong state, the force field binds the guest incorrectly, or the reference value was extracted from a different thermodynamic observable.

This distinction matters for materials discovery. A screening workflow can use a wrong adsorption score just as easily as a correct one. If the score is treated as a property of the material rather than as the output of a recipe, then a model may be optimised against a hidden mixture of structure error, force-field error, reference mismatch, and method mismatch. The result is not merely numerical noise. It is a false target.

The purpose of widom-atlas is to make that hidden recipe visible and testable. Each branch fixes one complete recipe before it is run. The branch then compares the resulting \(K_H\) and \(Q_{st}\), when those observables are defined, to one stated experimental reference. The outcome is not collapsed into a single success score. It is classified as a strict pass, a Tier-B physical-accuracy pass, a force-field disagreement, a reference-observable mismatch, a structural blocker, or an ensemble/method mismatch.

The distinction between these classes is essential. If a non-polarisable force field over-binds carbon dioxide at an open metal site, the answer is not to call Widom insertion wrong. If a cationic zeolite is represented by a pure-silica framework, the calculation is testing the wrong material. If a trapdoor zeolite opens during adsorption but the simulation samples only a rigid closed structure, the calculation and experiment are not evaluating the same partition function. These are different scientific failures, and treating all of them as ordinary pass/fail errors hides the physics.

Version 0.4 is therefore not presented as a general predictive adsorption tool, and the classical matrix is reported as method-plus-audited-examples rather than a fifteen-branch certification. Its locked result is narrow: two of fifteen verdict-affecting branches pass the strict criterion --- an all-silica MFI + Ar positive control \cite{Talu2001,Dunne1996Langmuir} and a UiO-66 + CO\(_2\) branch under genuine Universal Force Field (UFF) --- and five of fifteen pass the broader Tier-B physical-accuracy criterion, adding HKUST-1 + CO\(_2\) under genuine UFF, the Si-CHA + CO\(_2\) Bai 2013 branch compared with the Maghsoudi 2013 reference \cite{Bai2013JPCC,Maghsoudi2013Adsorption}, and the UiO-66 + CO\(_2\) Maia 2023 explicit-hydrogen variant \cite{Maia2023Crystals,Cmarik2012Langmuir}. The HKUST-1 and UiO-66 UFF passes were recovered only after the audit corrected a generator that had silently substituted guest and zeolite parameters under a false ``UFF'' framework label (Section~\ref{sec:result}). The Na-Rho + CO\(_2\) scalar branch is not counted as a scalar pass or fail because rigid closed-state Widom and open-state trapdoor adsorption are different observables \cite{Lozinska2012JACS,Verma2017ChemMater}.

The claim of this paper is consequently modest but useful. widom-atlas does not prove that Widom insertion with standard force fields is generally predictive. It proves that a Widom prediction can be made auditable branch by branch. In clean cases the recipe can reproduce experiment. In difficult cases the harness identifies whether the failure belongs to the force field, the reference, the structure, or the ensemble. That is the scientific object of the paper: not a universal adsorption calculator, but a governed way to decide whether a given Widom number should be trusted.

\textbf{Reader note on internal labels.} Codes such as \texttt{6c}, \texttt{3a}, \texttt{2a}, \texttt{4c}, and \texttt{3b\_EHq} are internal locked-recipe identifiers, not chemistry notation: each denotes one complete Widom-insertion recipe (a specific host structure, guest model, force field, charge model, mixing rule, electrostatics convention, simulation backend, sampling protocol, and reference observable). The ones used repeatedly below are branch 6c, the all-silica MFI + Ar positive control; branch 3a, UiO-66 + CO\(_2\) under genuine UFF; branch 2a, HKUST-1 + CO\(_2\) under genuine UFF; branch 4c, the Si-CHA + CO\(_2\) Bai 2013 control; and branch 3b\_EHq, the UiO-66 + CO\(_2\) Maia 2023 explicit-hydrogen, charged variant. ``Tier-A strict'' (\(|\Delta\log_{10}K_H|\le0.10\) and \(|\Delta Q_{st}|\le2.0\ \mathrm{kJ\,mol^{-1}}\)) and ``Tier-B physical-accuracy'' are the two pass bands defined in Section~\ref{sec:controls}. ``Native'' denotes this project's own independent Python Widom estimator --- a specific implementation, not a generic term --- used alongside the external Monte Carlo engines RASPA3 and RASPA2. The version labels are internal repository milestones: v0.4 is the locked classical audit matrix, v0.4.2 is the corrected rebuild and provenance layer for the audited examples, v0.5 is non-load-bearing consolidation or extension material, and v0.6 is the kUPS and MLIP governance companion (Section~\ref{sec:v06}).
